\chapter*{Resumen}
\markboth{Resumen}{}
\addcontentsline{toc}{chapter}{Resumen}

\ph{Resumen ejecutivo del TFG en español, normalmente entre 150 y 300 palabras.
Debe contener, en este orden:}

\begin{itemize}
  \item \ph{El contexto y problema que se aborda.}
  \item \ph{La propuesta o solución desarrollada (qué se ha hecho).}
  \item \ph{Las tecnologías o métodos clave utilizados.}
  \item \ph{Los resultados principales obtenidos, con 1--2 cifras concretas
  (p.~ej.\ cobertura de pruebas, tiempo de respuesta, número de usuarios o
  requisitos cubiertos): el tribunal lee el resumen antes que nada, y un
  resumen con números transmite un trabajo medido.}
  \item \ph{La conclusión o aportación principal.}
\end{itemize}

\ph{Redacta aquí el resumen en un único párrafo, sin viñetas, una vez completado el trabajo.}

\vspace{1.5em}
\noindent\textbf{Palabras clave:} \ph{5--10 términos separados por comas que
identifiquen los conceptos centrales (dominio, tecnologías, métodos, herramientas).}

\vspace{3em}

{\noindent\normalfont\Large\scshape Abstract}\\[-0.4em]
\noindent\rule{\textwidth}{0.4pt}
\addcontentsline{toc}{chapter}{Abstract}

\vspace{1em}

\ph{English translation of the abstract. Same structure and length as the
Spanish version, including the same 1--2 concrete figures for the main
results.}

\vspace{1.5em}
\noindent\textbf{Keywords:} \ph{5--10 comma-separated terms in English.}
